\chapter{如何說話，讓人想聽}

這場講座篇幅不長，卻以異乎尋常的精密結構建成。Treasure 先把賭注拉高：人的聲音是我們每一個人都在演奏的樂器，也許是世上最有力量的聲音。接著，他又把這份恢弘收束為一個實際的難題。若這種媒介如此強大，為何人們明明開口說話，卻常常無人傾聽？他的回答按著一條刻意安排的次序展開：先診斷那些摧毀可聽性的習慣，再提出一個正面的倫理基礎，然後給出一組聲音變數的工具箱，接著是一套熱身程序，最後才擴展到一個更大的願景，去想像一個自覺對待聲音的世界會是什麼樣子。

\section{聲音與問題}

開場的主張刻意取得很大。聲音可以引發一場戰爭；它也可以說出 ``I love you.'' Treasure 希望我們首先感受到，言語不是裝飾，它是有後果的。只有在這樣的放大之後，他才提出這場演講的主導問題：為什麼人們那麼常聽不進別人的話，而言語又該如何變得足夠有力，甚至足以改變世界？

若用一種簡潔的編者速記，我們可以把這個悖論概括如下：
\begin{equation}
\text{強大的媒介} \not\Rightarrow \text{有效的溝通}.
\end{equation}
聲音蘊藏巨大的潛能，但潛能本身並不保證會被接收。

Treasure 接下來這一步至關重要。他並不是先給出一個讓聲音變好聽的小技巧，而是先問：究竟是什麼先把傾聽堵住了？因此，整場講座是先診斷，再處方。

\subsection{問題與回答}

\paragraph{問題。} 為什麼我們說話時，人們卻聽不進去？

\paragraph{回答。} Treasure 的答案有兩層。首先，在技巧還沒登場之前，就已經有一系列反覆出現的習慣，使我們變得難以被傾聽。其次，即便我們的意圖是好的，表達方式依然重要：同樣的字句，會因為說出的方式不同，而帶著不同的力量、權威與意義抵達聽者。因此，這場講座便從習慣，走向基礎，再走向聲音機制。

\section{說話的七宗罪}

Treasure 的第一個回答完全是負面的。他列出一份 ``seven deadly sins of speaking''，並不是要建立一套完整的道德體系，而是要對那些腐蝕可聽性的習慣做出一種實際的分類。它們的順序本身就是講座節奏的一部分，因此值得明確保留下來：
\begin{equation}
S_7=
\bigl(
\text{閒話},
\text{論斷},
\text{負面},
\text{抱怨},
\text{藉口},
\text{誇飾},
\text{武斷}
\bigr).
\end{equation}

\begin{remark}
這個記號是編者所加，其目的只是為了在紙面上保留講座這套有次序的清單。
\end{remark}

現在我們可以按著講座本身的步伐，一一走過這份名單。

\begin{enumerate}
\item \textbf{閒話。} 講別人的壞話，而那人又不在場，會立刻削弱信任。若一位講者會拿別人來說嘴，我們自然會推想：同一個人日後也會在別處說我們。
\item \textbf{論斷。} 當人一面感到自己正被審視、被判定不足，一面還要去聽對方說話，這是很困難的。愛論斷的講者，會把聽者逼進自我防衛之中。
\item \textbf{負面。} Treasure 用關於十月一日多麼糟糕的那句話來示範。重點不只是滑稽而已。持續的負面會收縮聽者的世界，使說話變得沉重。
\item \textbf{抱怨。} 它被當作另一種負面形式提出，卻仍值得獨立站出來，因為它太常見了。Treasure 稱之為病毒式的痛苦：它傳播的不是光亮，而是不快。
\item \textbf{藉口。} 這裡的典型人物，是那種到處推責任的人，把一切都往別人身上丟。這種說話之所以讓人覺得閃躲，正因為它從不讓責任真正落定。
\item \textbf{誇飾。} 他所說的 embroidery，意思就是誇張。語言失去了刻度。若任何事都被稱作 awesome，那麼語言就再也無法標示真正的敬畏。而誇張若推到極致，也就變成了說謊。
\item \textbf{武斷。} 武斷，就是把事實與意見混為一談。一旦兩者糾纏不清，聽者就不再是被邀請進入思考，而是被一連串冒充真理的斷言迎面擊打。
\end{enumerate}

這裡重要的是一種累積性的結構。Treasure 並不只是叫我們更和善一點。他是在逐一展示：講者與聽者之間的通道，是如何被一種又一種習慣損壞的。到了清單的末尾，診斷便變得銳利了：人們之所以不聽，往往是因為我們自己早已把傾聽變成了一件困難的事。

也正因此，講座在此刻需要一個真正的轉折。診斷之後，便要求一個基礎。

\section{HAIL：正面的基礎}

Treasure 現在提出前一節已經準備好的問題：有沒有一種正面的方式來理解這件事？有。答案是另一個有次序的結構，被壓縮成助記詞 \(\mathrm{HAIL}\)：
\begin{align}
H &\mapsto \text{誠實},\\
A &\mapsto \text{真實},\\
I &\mapsto \text{正直},\\
L &\mapsto \text{善意}.
\end{align}

這個詞之所以重要，有兩層意思。它容易記住；同時，它也帶有問候或熱烈致意的意味。暗示正在於此：若一個人的話語是從這四種基礎上說出的，它也會以那樣的方式被接住。

\subsection{問題與回答}

\paragraph{問題。} 取代七宗罪的正面基礎是什麼？

\paragraph{回答。} Treasure 的回答是，有力的說話必須是真實的、真誠的、可靠的，也是懷著善意的。誠實使那條線不致歪斜。真實意味著我們不是戴著面具在說話。正直意味著我們言行一致。善意則意味著我們希望聽者好。這些不是技巧之外的裝飾，而是技巧得以站立其上的地基。

這一節裡最有意思的局部張力，出現在誠實與善意的關係中。Treasure 並不主張 ``absolute honesty''，彷彿只要絕對誠實就已足夠。他立刻給出一個反例：毫無必要的直白傷人。若用緊湊的形式來寫：
\begin{equation}
\text{絕對誠實} \not\Rightarrow \text{好的說話},
\qquad
\text{以善意調和的誠實} \Rightarrow \text{好的說話}.
\end{equation}
善意並不取消真實；它規約真實的使用方式。

接著，他又補上一個較柔和的主張：
\begin{equation}
\text{善意} \;\Rightarrow\; \text{較少傾向於論斷}.
\end{equation}
這不是一條定理，也不該被當作定理來讀。它指的是一種質性的相斥關係：若我們真的在祝願某人好，那麼用論斷的姿態去面對那個人，就會變得更困難。

Treasure 以一個果斷的轉折結束這一節：以上說的是我們說什麼。接下來還要問：我們如何去說。

\section{聲音的工具箱}

講座到了這裡，開始從倫理基礎轉向聲音機制。借用 Treasure 的說法，聲音成了一個神奇的工具箱，而大多數人從未真正打開過它。若是為了做筆記，把這些工具收攏成一個整體會很有幫助：
\begin{equation}
V_{\text{tool}}=
\bigl(
\text{音域位置},
\text{音色},
\text{語調節奏},
\text{速度},
\text{沉默},
\text{音高},
\text{音量}
\bigr).
\end{equation}

重點不是 Treasure 真的寫下了這樣一條公式，他沒有。重點在於，講座本身現在就是按著這樣一種方式運作，彷彿說話是一個可控制的系統，而其中有一小組主要變數。

\paragraph{音域位置。} Treasure 首先對比聲音被放置的位置：高高地放在鼻腔，更居中的喉部，以及更低的胸腔。他並不是在開一堂正式的聲學課，而是在描畫一種說話的實用現象學。他說，如果我們想要重量感，往往就會往下走。這便直接引向他那個帶有修辭色彩的經驗性主張：
\begin{equation}
\text{較低的音域位置} \;\Rightarrow\; \text{更強的力量感與權威感}.
\end{equation}
這也就是為什麼他說，我們會投票給那些聲音較低的政治人物。重點不在數值測量，而在感受到的力量。

\paragraph{音色。} 音色，是聲音給人的質地感。Treasure 把理想的聲音描述為豐富、圓潤、溫暖，``like hot chocolate.'' 隱藏在這個比喻中的技術重點很重要：音色是可以訓練的。呼吸、姿勢與練習，都能改變它。

\paragraph{語調節奏。} 對 Treasure 而言，prosody 是對話中意義的第一通道。它是那種帶著起伏的元語言，使話語獲得力量與方向。這就是為什麼單調的說法難以入耳：它拿掉了意義最主要的載體之一。

\subsection{問題與回答}

\paragraph{問題。} 為什麼同樣的字句，在不同說法下會有不同意思？

\paragraph{回答。} 因為意義並不只棲居在詞語串列本身之中。Treasure 那些關於句尾反覆上揚的語調，以及 ``Where did you leave my keys?'' 這句話的例子，都在說明：即使字詞固定，輪廓仍會改變這個言語行動本身。

我們可以把這個觀念以最簡約的形式寫下。令
\begin{align}
u &= \text{``Where did you leave my keys''},\\
m_1 &= \mathcal{M}(u,\pi_1),\\
m_2 &= \mathcal{M}(u,\pi_2),
\end{align}
其中 \(\pi_1\) 與 \(\pi_2\) 是不同的語調或音高輪廓，而 \(\mathcal{M}\) 是聽者所感知到的意義。那麼 Treasure 的要點不過是
\begin{equation}
m_1 \neq m_2.
\end{equation}

\begin{remark}
這同樣是編者的記號。它只是捕捉這個示範的邏輯：字詞相同，輪廓不同，意義也就不同。
\end{remark}

這正是為什麼反覆把句尾講成疑問上揚會成為問題。當每一句話都上揚得像一個問句時，陳述句就失去了它作為陳述的力量。同一種輪廓，無法承載全部本來打算傳遞的意義範圍。

\paragraph{速度與沉默。} 接著，Treasure 透過對比來示範速度。人可以說得很快，以製造興奮；也可以慢下來，以形成強調。而在這個尺度的盡頭，便是沉默。沉默並不是未能繼續說話；它本身就是一種溝通工具：
\begin{equation}
\text{快速} \Rightarrow \text{興奮},
\qquad
\text{緩慢} \Rightarrow \text{強調},
\qquad
\text{沉默} \Rightarrow \text{有力的停頓}.
\end{equation}
這也就是為什麼，一場演講不必被 ``ums'' 和 ``ahs'' 填滿。沉默本身可以承載重量。

\paragraph{音高。} 音高常常與速度一起作用，用來標示亢奮程度，但 Treasure 也示範了：它可以單獨發揮作用。同一句話的兩種說法之所以不同，是因為旋律輪廓不同。字詞保持不變，意圖卻改變了。

\paragraph{音量。} 最後一件工具是音量。大聲可以製造命令感與興奮感。小聲則能讓整個房間向前傾身。Treasure 警告的，是那種持續不斷的放送，他把它稱作 ``sodcasting''：一種把聲音粗暴地施加給他人的做法。

整個這一節完成了一次真正的技術性轉移。我們不再只是被告知要 ``好好說話''；我們看見了一組可以被刻意調節的變數，而每一個變數都會改變聽者接收話語的方式。

\section{準備這件樂器}

下一個轉折是動機性的。Treasure 要我們想一想那些說話真正重要的時刻：公開演講、提出企劃、要求加薪、婚禮致詞。在這樣的時刻，他說，我們有責任為自己做好準備。此時，那個比喻再次收緊：聲音不只是一個工具箱，它也是一具引擎，而任何引擎若未經暖機，都不會運作得好。

\subsection{問題與回答}

\paragraph{問題。} 為什麼在說話之前要先暖聲？

\paragraph{回答。} 因為重要的說話值得好好照料這件樂器。若聲音是承載整個溝通行動的引擎，那麼我們就不該在冰冷狀態下啟動它。Treasure 並不把這件事停留在比喻層面，而是把它轉成一套具體的演出前程序。

若為求簡潔參照，我們可以寫成
\begin{equation}
W_6=
\bigl(
\text{呼吸與嘆氣},
\text{嘴唇敲擊音},
\text{嘴唇顫音},
\text{舌頭 } la\text{ 練習},
\text{捲舌 } r,
\text{警笛音 } (we \to or)
\bigr),
\end{equation}
但真正的內容在於這套有次序的程序本身：
\begin{enumerate}
\item 抬起雙臂，深深吸氣，再把氣嘆出來。
\item 用重複的 ``bop'' 聲來暖開嘴唇。
\item 接著做嘴唇顫音，就像孩子會發出的那種嗡嗡聲。
\item 用誇張的 ``la, la, la'' 發音把舌頭喚醒。
\item 捲一個 \(r\) 音；Treasure 很生動地稱它為 ``champagne for the tongue.''
\item 最後做警笛音，從高音的 ``we'' 滑到低音的 ``or.''
\end{enumerate}

這套程序之所以重要，是因為它構成了整場講座第一個完全實用的回報。Treasure 已經從惡習走到基礎，再走到變數，而現在則走到程序。聽眾不只是被告知什麼重要；他們還被要求去預演它。

接著，他用一句簡短的命令式尾聲收束整個段落：下次你要說話時，事先把這些做一遍。這句話應當在筆記中清楚留存，因為它標誌著：論述已經在此轉化為習慣。

\section{結尾的願景：有意識地說、有意識地聽、有意識地創造聲音}

在結尾處，Treasure 說，他想把整場討論放回更大的脈絡中。這個動作很有代表性：在一連串緊湊的清單與局部示範之後，他把畫面拉開，開始問，這些做法究竟會建成一個什麼樣的世界。

他對當下狀態的診斷，可以簡潔地寫成
\begin{equation}
\text{當前世界}
=
\bigl(
\text{糟糕的說話},
\text{糟糕的傾聽},
\text{噪音},
\text{不良的聲學}
\bigr).
\end{equation}
這不只是個人的缺陷，也是環境性的缺陷。人說得不好；聽者也不聽；連聲學環境本身都在妨礙理解。

而他要我們想像的反事實世界，則正好相反：
\begin{equation}
\text{理想世界}
=
\bigl(
\text{有力的說話},
\text{有意識的傾聽},
\text{有意識地創造聲音},
\text{適配目的的聲學}
\bigr).
\end{equation}

因此，講座並不是以一堂私人的自我改進課結束，而是以一種公共的聲學視野作結。如果我們有意識地創造聲音、有意識地接收聲音，也有意識地為聲音設計我們的環境，會發生什麼事？Treasure 的回答是：那樣的世界將會聽起來很美，而理解將成為常態。

這最後一次擴展，對整場講座的節奏至關重要。最初那個關於人們為何不聽的問題，最後竟變成了一個更大的問題：我們究竟正用自己的說話與傾聽習慣，建造出一個什麼樣的人類世界？

\section{總結}

Treasure 的講座，是由一系列彼此嵌套、且每一步都由前一步所驅動的結構組成的。我們從那個強大聲音卻常常失效的悖論開始，接著透過七宗罪來診斷這種失敗，再藉由 HAIL 重建其道德地基。我們打開聲音的工具箱，看見表達並非偶然，而是由可調節的變數構成。接著，我們用六個熱身動作來準備這件樂器。最後，整個探問又被擴展到一個有意識對待聲音的世界之中。

因此，這一章不應被讀成一袋零散的說話小技巧，而應被讀成一套小而完整、卻結構分明的溝通理論。言語之所以重要，是因為它有力量。它會因為可辨認的理由而失效。它可以建立在清楚的基礎上重新被鍛造。它的力量依賴於可控制的變數。而當我們學會有意識地運用這些變數，收穫的不只是修辭上的有效性，更是一種更大規模的理解可能。